% 3.1 布拉维格子
\section{布拉维格子}

每个空间群 $\mathbf{G}$ 都有一组纯平移 $\{E \,|\, \mathbf{t}\}$，它们构成空间群 $\mathbf{G}$ 的一个不变子群 $\mathbf{T}$。这就是晶体所依托的布拉维格子的平移对称操作群；$\mathbf{t}$ 由式 (1.5.1) 给出：

\begin{equation*}
\mathbf{t} = n_{1}\mathbf{t}_{1} + n_{2}\mathbf{t}_{2} + n_{3}\mathbf{t}_{3},
\tag{1.5.1}
\end{equation*}

其中 $n_{1}$、$n_{2}$、$n_{3}$ 是整数，$\mathbf{t}_{1}$、$\mathbf{t}_{2}$、$\mathbf{t}_{3}$ 是布拉维格子的基本平移。如 1.5 节所述，由向量 $\mathbf{t}$ 确定的点称为格点；建立在向量 $\mathbf{t}_{1}$、$\mathbf{t}_{2}$、$\mathbf{t}_{3}$ 上的平行六面体称为基本晶胞。基本晶胞与惯用晶胞的区别已在 1.5 节提及并在图 1.8 中说明。整个格子由大量晶胞组成，相邻晶胞之间相差格子的一个向量 $\mathbf{t}$。严格地说，布拉维格子是无穷延伸的。固然可以假定真实晶体无穷延伸，然后发展其空间群及平移子群的表示理论；但实践中通常假定晶体有限，并施加周期性边界条件：设沿 $\mathbf{t}_{i}$（$i = 1, 2, 3$）方向有 $N_{i}$ 个晶胞，$N_{i}$ 是很大的数。为使 $\mathbf{G}$ 仍是群，采用 Born--von Kármán 周期性条件（Born and von Kármán 1913*a*, *b*）：若 $l_{1}$、$l_{2}$、$l_{3}$ 是整数，则

\begin{equation*}
\{E \,|\, l_{1}N_{1}\mathbf{t}_{1} + l_{2}N_{2}\mathbf{t}_{2} + l_{3}N_{3}\mathbf{t}_{3}\} = \{E \,|\, \mathbf{O}\}.
\tag{3.1.1}
\end{equation*}

这样做的好处是空间群成为有限群，从而可以应用有限群的理论。习惯上随后令 $N_{i} \rightarrow \infty$；这样做的理由是，由此得到的是一切具有物理意义的结果。直接处理无穷群并无不可，只是对无穷群必须重新表述不可约表示的许多定义和定理。

在完整描述各种可能的格子之前，先作几点一般性说明。第一，所有初基晶胞体积相同，均为

\begin{equation*}
V = \mathbf{t}_{1} \cdot (\mathbf{t}_{2} \times \mathbf{t}_{3}).
\tag{3.1.2}
\end{equation*}

第二，建立在布拉维格子任意三个独立格矢上的平行六面体，其体积不小于 $V$。第三，表征格子的基本矢量 $\mathbf{t}_{1}$、$\mathbf{t}_{2}$、$\mathbf{t}_{3}$ 常可有多种取法；我们将作出一个明确的取法（见表 3.1），此后即认定该取法。

1.5 节说过可能有 14 种不同的布拉维格子，并已在图 1.7 中绘出。同一晶系的两种布拉维格子，若能通过一个不经过更低对称晶系的连续形变互相变换，则没有区别。例如面心立方布拉维格子 $\Gamma_{c}^{f}$ 不能连续形变为体心立方布拉维格子 $\Gamma_{c}^{v}$，除非经过三方布拉维格子 $\Gamma_{rh}$。表 3.1 列出全部 14 种布拉维格子及我们所用的基本平移和相应晶胞的体积。晶体结构或空间群的\textit{类型}就是它所依托的布拉维格子，通常用表 3.1 第 1 列的符号标记。对单斜晶系，表 3.1 采用《X 射线晶体学国际表》（Henry and Lonsdale 1965）给出的两种定向中的第一种；$a$、$b$、$c$ 是惯用晶胞相邻棱上的向量，在《国际表》中，除布拉维格子对称性所要求的以外，对 $a$、$b$、$c$ 的相对大小没有其他限制。

表 3.1 的注（译文见下）逐列说明了各列含义；特别地，第 3 列给出相对右手直角坐标系 $Oxyz$ 表示的布拉维格子基本平移：例如 $\mathbf{t}_{1} = (p, q, r)$ 意为 $\mathbf{t}_{1} = p\mathbf{i} + q\mathbf{j} + r\mathbf{k}$，其中 $\mathbf{i}$、$\mathbf{j}$、$\mathbf{k}$ 是 $x$、$y$、$z$ 方向的单位向量；第 4 列以惯用晶胞的棱长 $a$、$b$、$c$ 给出初基晶胞体积（按式 (3.1.2)）。注 (v) 指出，《国际表》对某些空间群（表 3.7 的编号 38--41）对正交底心布拉维格子 $\Gamma_{o}^{b}$ 使用了另一种取向（记作 A），本书不使用 A 格子。注 (vii) 说明单斜晶系采用第一种定向，即把二次对称轴取在 $z$ 轴上。注 (viii) 给出 Belov（1964）汇集的 14 种布拉维格子的其他记号（Schönflies 记号、国际记号、Wilson 记号与 Pearson 记号的对照表）。尽管 14 种布拉维格子如此重要，目前还没有一种被普遍接受的、能简洁而清晰地描述每种格子的记号；《国际表》使用的公认字母 P、I、F、B、C、R（表 3.1）只表明格子是初基的或以某种方式有心。Schönflies 符号中的字母 F 是多余的。Wilson 建议对每种格子用一个专门字母，Pearson 建议用双字母符号；这些备选方案由 Belov（1964）汇总为表 3.1 的注 (viii)。

若低对称布拉维格子的棱长 $a$、$b$、$c$ 之间存在某些特定关系，它就可能升入更高对称的布拉维格子。例如 $a = c/\sqrt{2}$ 时 $\Gamma_{rh}$ 成为 $\Gamma_{c}^{f}$；$a = \sqrt{2}\,c$ 时 $\Gamma_{rh}$ 成为 $\Gamma_{c}$；$a = 2\sqrt{2}\,c$ 时 $\Gamma_{rh}$ 成为 $\Gamma_{c}^{v}$——即所有立方格子都是三方格子的特例。又如 $a = c$ 时 $\Gamma_{q}$ 成为 $\Gamma_{c}$、$\Gamma_{q}^{v}$ 成为 $\Gamma_{c}^{v}$；$c = \sqrt{3}\,b$ 时 $\Gamma_{o}^{b}$ 成为 $\Gamma_{h}$。其他此类关系不难由表 3.1 推出。

\begin{figure}[H]
\centering
\includegraphics[width=0.92\textwidth]{c3_t31_p1.png}
\end{figure}
\begin{figure}[H]
\centering
\includegraphics[width=0.92\textwidth]{c3_t31_p2.png}
\caption*{\textbf{表 3.1}\ 14 种布拉维格子：格子类型符号、Schönflies 记号 $\Gamma$、国际记号、基本平移（相对 $Oxyz$）与初基晶胞体积（原书扫描占位，含注 (i)--(viii)，见正文）。}
\end{figure}

七个晶系的对称元素见表 1.2。了解晶系的点群操作对格子平移的作用是有益的，这些作用列于表 3.2。
